\section{Large covariance matrices: shrinkage and random-matrix noise}

\begin{qfmethod}[Estimating a large covariance matrix: why the sample matrix fails when assets are many relative to observations, and linear and nonlinear shrinkage]\label{meth:portfolio-covariance-estimation}
\qffield{Statement} Ledoit and Wolf's review starts from the case ``when the number of variables is of the same magnitude as the number of observations'', where ``the sample covariance matrix certainly will not do'' \citep[UZH ECON WP~323, version of February 2020, abstract, p.~1]{wide-ledoit-wolf-review}. \emph{Linear shrinkage} to the identity solves $\min_{\rho_1,\rho_2}E\|\hat\Sigma-\Sigma\|_F^2$ over $\hat\Sigma=\rho_1I+\rho_2S$, with the scaled Frobenius norm $\|A\|_F^2=\operatorname{Tr}(A'A)/N$ (eqs.~(2.1)--(2.5), pp.~4--5). The solution is
\[ \Sigma^*=\gamma^*\mu I+(1-\gamma^*)S,\qquad \gamma^*=\frac{\beta^2}{\delta^2}, \]
with $\mu=\langle\Sigma,I\rangle$ the average variance, $\alpha^2=\|\Sigma-\mu I\|_F^2$, $\beta^2=E\|S-\Sigma\|_F^2$ and $\delta^2=E\|S-\mu I\|_F^2=\alpha^2+\beta^2$ (eqs.~(2.6)--(2.7), p.~5). It is a convex combination, positive definite and invertible ``even when $N>T$''. The optimal intensity depends on unknown population quantities, so a feasible estimator is derived under large-dimensional asymptotics, with $N/T\to c$ (sec.~2.2, p.~6).
\begin{itemize}
\item \emph{Other targets.} Linear shrinkage can also target the single-factor matrix or the constant-correlation matrix.
\item \emph{Their own comparison.} In their 2004 back-test of a tracking-error-constrained manager, the single-factor target, the constant-correlation target and a five-principal-component factor model all ``dominated the sample covariance matrix, but there was no clear winner'': constant correlation was best for $N\le100$ and the single factor for $N\ge225$ (sec.~3.3, p.~14).
\item \emph{Nonlinear shrinkage} keeps the sample eigenvectors and shrinks each eigenvalue differently. It ``can deliver another level of performance improvement, especially if overlaid with'' time-varying co-volatility or factor models (abstract; secs.~4--6).
\end{itemize}

\qffield{Notation} $N$ assets, $T$ observations, $S$ sample covariance, $c=N/T$ the concentration; the norm divides by $N$, unlike the Frobenius norm.

\qffield{Conditions} Section~2 assumes i.i.d., mean-zero observations; demeaned data and dynamic (DCC) and factor extensions are treated later (Remark~2.1; secs.~5--6). Optimality is in Frobenius loss, which is not the same as the out-of-sample variance of an optimised portfolio, though the review's applications are mostly Markowitz problems.

\qffield{Practice} Own simulation with a known one-factor population of 100 assets and 120 observations:
\begin{itemize}
\item $\alpha^2+\beta^2=\delta^2$ holds (0.2035 against 0.2031);
\item the oracle intensity is 0.045, and the Frobenius loss falls from 0.00910 to 0.00871;
\item the true variance of the global minimum-variance portfolio falls from 0.0252 using $S$ to 0.0089 using $\Sigma^*$, against an optimum of 0.0044.
\end{itemize}
The small loss gain and the large portfolio gain illustrate that the two criteria differ. Narang's grouping of similar instruments (\hyperref[mod:factor-risk-model]{factor risk model}) is a structured target of this kind, without a rule for the intensity.

\begin{itemize}
\item \emph{The constant-correlation target and its feasible intensity} (\nolinkurl{s6-honey}, 2003 manuscript). The target $F$ keeps the sample variances and sets every covariance to $\bar r\sqrt{s_{ii}s_{jj}}$, with $\bar r$ the average sample correlation (App.~A, eq.~(3), p.~12). Under Frobenius loss the optimal weight on $F$ behaves like $\kappa/T$ with $\kappa=(\pi-\rho)/\gamma$: $\pi$ sums the asymptotic variances of the $\sqrt T s_{ij}$, $\rho$ the asymptotic covariances of $\sqrt Tf_{ij}$ with $\sqrt Ts_{ij}$, and $\gamma=\sum_{ij}(\phi_{ij}-\sigma_{ij})^2$ measures the target's misspecification. Each is replaced by a consistent sample estimate, and the intensity is $\hat\delta^*=\max\{0,\min\{\hat\kappa/T,1\}\}$ (App.~B, eqs.~(4)--(5), pp.~12--15) \citep[App.~A--B, pp.~12--15]{s6-honey}.

Numerically, with 50 assets and 60 observations the shrunk matrix has the lower Frobenius loss; the intensity goes to one when the population really has a constant correlation (the target is then right) and is about 0.3 when correlations come in blocks; with more assets than observations the shrunk matrix is invertible and the sample matrix is not.
\item \emph{POET} (\nolinkurl{wide-fan-liao-mincheva-poet}, arXiv 1201.0175). Keep the first $K$ principal components of the sample covariance and threshold the rest, the ``principal orthogonal complement'': off-diagonal entries $\hat r_{ij}$ of the remainder are set to $s_{ij}(\hat r_{ij})I(|\hat r_{ij}|\ge\tau_{ij})$ with an adaptive threshold $\tau_{ij}=\tau(\hat r_{ii}\hat r_{jj})^{1/2}$, so $\tau=0$ gives the sample covariance and $\tau=1$ a strict factor model (sec.~2, eqs.~(2.4)--(2.7), printed pp.~9--10). It assumes \emph{pervasive} factors (the eigenvalues of $p^{-1}B'B$ bounded away from zero, so the first $K$ eigenvalues grow with the number of assets $p$) and a sparse residual covariance \citep[sec.~2, printed pp.~6--10]{wide-fan-liao-mincheva-poet}.
\item \emph{Dynamic correlation} (Engle and Sheppard's account of DCC, NBER WP~8554, 2001; Engle's 2002 \emph{JBES} paper is \nolinkurl{wide-engle-dcc}). Each series gets its own univariate GARCH; the standardised residuals $\epsilon_t$ drive $Q_t=(1-\sum\alpha_m-\sum\beta_n)\bar Q+\sum\alpha_m\epsilon_{t-m}\epsilon'_{t-m}+\sum\beta_nQ_{t-n}$, with $\bar Q$ the unconditional covariance of the standardised residuals, and the correlation matrix is $R_t=Q_t^{*-1}Q_tQ_t^{*-1}$ with $Q_t^*$ the square roots of $Q_t$'s diagonal; $\alpha_m,\beta_n\ge0$ with $\sum\alpha_m+\sum\beta_n<1$ and a positive-definite $\bar R$ keep $R_t$ a correlation matrix (eqs.~(1)--(2) and Proposition~2, printed pp.~5--6) \citep[secs.~2--3, printed pp.~5--6]{engle-sheppard-2001-dcc}. Estimation is in two steps, the variances first and then the correlation parameters.
\item \emph{The published review} (\emph{Journal of Financial Econometrics} 20(1), 2022). As a rule of thumb, ``one generally needs both $N$ and $T$ to be above 50 for nonlinear shrinkage to outperform linear shrinkage'' when the optimal transformation is markedly nonlinear, and then ``there is basically nothing to lose but potentially a lot to gain by upgrading'' (sec.~3.3, pp.~205--206). For Markowitz problems with time-varying co-volatility they propose DCC-NL: univariate GARCH(1,1) to devolatilise, nonlinear shrinkage of the unconditional correlation matrix, and composite likelihood for the dynamics; RiskMetrics 2006 ``performed poorly and cannot be recommended for the purpose of Markowitz's portfolio selection'' (sec.~4, pp.~208--210, n.~13) \citep[2022 journal version, secs.~3.3--4, pp.~205--210]{wide-ledoit-wolf-review}. The working paper read earlier and the journal version agree on the passages compared (abstract and the linear-shrinkage solution).
\end{itemize}

\qffield{Pitfalls} The nonlinear estimator itself was not transcribed. For how the loadings, factor covariance and residual variances of a factor model are estimated, see Connor, Goldberg and Korajczyk (\nolinkurl{wide-connor-goldberg-korajczyk}). \emph{Time-varying estimators}. Boyd et al.\ list EWMA, Engle's DCC and iterated EWMA as approaches that ``work well in practice'', with or without a factor model. They give no comparison and refer to Johansson et al.\ (2023) for detail \citep[arXiv v1, sec.~3.4, pp.~21--22]{wide-boyd-markowitz-seventy}. In their backtest the estimator is $\hat\Sigma_t=\alpha_t\sum_{\tau=1}^t\beta^{t-\tau}r_\tau r_\tau^\top$ with $\alpha_t=(1-\beta)/(1-\beta^t)$. It uses the second moment, the mean being negligible, with a half-life of 125 trading days, $\beta\approx0.994$ ($0.5^{1/125}=0.9945$). The authors note that the half-life does not change their results qualitatively (pp.~31--32). Interpretation: the normalisation makes the weights sum to one in short samples, which a fixed $1-\beta$ does not. The univariate GARCH and EWMA comparison is in \hyperref[mod:garch-volatility-forecast]{GARCH volatility forecasts}.

\qffield{See also} \hyperref[mod:factor-risk-model]{factor risk model}; \hyperref[mod:mean-variance-portfolio-optimisation]{mean--variance optimisation}.

\qffield{Further reading} \nolinkurl{wide-connor-goldberg-korajczyk}.

\end{qfmethod}

\begin{qfevidence}[How much of an empirical correlation matrix is noise: the Marchenko--Pastur band against 406 S\&P~500 stocks]\label{ev:correlation-matrix-noise}
\qffield{Statement} For a $T\times N$ matrix of independent entries of variance $\sigma^2$, with $Q=T/N\ge1$ fixed as both grow, the eigenvalue density of the correlation matrix is $\rho(\lambda)=\frac{Q}{2\pi\sigma^2}\frac{\sqrt{(\lambda_{\max}-\lambda)(\lambda-\lambda_{\min})}}{\lambda}$ on $[\lambda_{\min},\lambda_{\max}]$, with $\lambda_{\max,\min}=\sigma^2(1+1/Q\pm2\sqrt{1/Q})$ \citep[arXiv cond-mat/9810255v1, eq.~(3), p.~1]{wide-laloux-noise-dressing}. Laloux, Cizeau, Bouchaud and Potters compare it with the correlation matrix of $N=406$ S\&P~500 stocks over $T=1309$ days, 1991--96 ($Q=3.22$). The largest eigenvalue, the market mode, is ``25 times larger than the predicted $\lambda_{\max}$''. Taking the rest as noise and fitting $\sigma^2$, the best fit ($\sigma^2=0.74$) ``accounts quite satisfactorily for 94\% of the spectrum, while the 6\% highest eigenvalues still exceed the theoretical upper edge'' (pp.~1--2).

\qffield{Notation} $N$ assets, $T$ observations, $Q=T/N$; returns demeaned and scaled to unit volatility before the correlation matrix is formed.

\qffield{Conditions} Evidence on 1991--1996 data. The band is an $N\to\infty$ limit for independent, identically distributed noise; for finite $N$ the edges blur. Heavy tails and volatility clustering in returns also move the edges, which the authors say a smaller effective $Q$ could absorb (p.~2). Eigenvalues inside the band are not proved to be noise, only indistinguishable from it on this test.

\qffield{Practice} With pure noise at the paper's $N$ and $T$, the simulated eigenvalues fall within the band $[0.196,2.424]$ up to the finite-sample blur. \emph{A slip in the source} (own reading): the text derives the residual variance as ``$\sigma^2=1-\lambda_{\max}/N=0.85$''. With $\lambda_{\max}=2.43$ that would give 0.994; with the market eigenvalue $\lambda_1=25\lambda_{\max}\approx60.6$, $1-\lambda_1/406=0.851$, so the formula must mean $\lambda_1$. The practical use is in \hyperref[meth:portfolio-covariance-estimation]{covariance estimation}: clipping or shrinking the eigenvalues inside the band, which Bun, Bouchaud and Potters review (\nolinkurl{wide-bun-bouchaud-potters-cleaning}).

\qffield{Pitfalls} The sample is 1991--96 daily data for one index; the 94/6 split is a fit with $\sigma^2$ chosen to match, not a test. Locators refer to the arXiv preprint of the letter (\emph{Physical Review Letters} 83, 1999).

\qffield{See also} \hyperref[meth:portfolio-covariance-estimation]{covariance estimation}; \hyperref[mod:pca-risk-distribution-weights]{PCA risk-distribution weights}.

\qffield{Further reading} \nolinkurl{wide-bun-bouchaud-potters-cleaning}; \nolinkurl{wide-vershynin-hdp}.

\end{qfevidence}
